
\documentclass[10pt]{article}
\usepackage[margin=0.82in]{geometry}
\usepackage{xcolor}
\usepackage{enumitem}
\usepackage{titlesec}
\usepackage{setspace}
\usepackage{hyperref}
\setstretch{1.04}
\setlength{\parindent}{0pt}
\setlength{\parskip}{0.45em}
\setlist[itemize]{leftmargin=1.4em,itemsep=0.15em,topsep=0.2em}
\titleformat{\section}{\Large\bfseries}{\thesection}{0.5em}{}
\titleformat{\subsection}{\large\bfseries}{\thesubsection}{0.5em}{}
\newcommand{\coachcomment}[2]{%
  \par\noindent\fcolorbox{orange!70!black}{yellow!20}{\parbox{0.965\linewidth}{\textbf{Highlighted coach/review comment - #1:} #2}}\par
}
\newcommand{\tableblock}[2]{%
  \par\noindent\fcolorbox{black!25}{gray!6}{\parbox{0.965\linewidth}{\textbf{#1}\par #2}}\par
}
\begin{document}
\begin{center}
{\Large\textbf{Reviewed Draft Excerpt With Coach Comments Highlighted}}\\[0.3em]
EDCI 59100-016: Literature Review\\
Kamrul Md Kamruzzaman\\
\end{center}
\noindent\textbf{Scope of this PDF:} This excerpt contains the draft through Chapter 3 only. Chapter 4, Chapter 5, and the bibliography/reference section are omitted as requested. Word comments and advisor notes are shown in highlighted boxes.

\section*{Introduction}
Construction job sites are dynamic environments involving various workers, equipment, and machines operating simultaneously, resulting in managing such job sites a complex and challenging task (Alaloul et al. 2022). Compared with some other engineering sectors such as automotive and aerospace, construction is heavily labor intensive, has a vast and diverse supply chain, and is dominated heavily by single point and ad-hoc solutions, which results in considering every project is different (Davila Delgado and Oyedele 2021). Because construction sites have these characteristics, standard industrial monitoring systems are often not practical for use in the construction industry (Alaloul et al. 2022).

Due to this profound complexity, regular safety inspections are needed to keep the work environment hazard-free (Asadzadeh et al. 2020). Additionally, managers must monitor and track activities and ensure the project meets the expected production rates (Alaloul et al. 2022). However, traditional manual observation and reporting methods, such as physical site walk-throughs for hazard identification or the manual logging of Daily Construction Reports (DCRs) heavily rely on subjective and error-prone manual visual inspection procedures (Samsami 2024). Furthermore, these manual methods are time-consuming, labor-intensive, and result in a slow information flow (Alaloul et al. 2022).

The absence of real-time information resulting from these manual limitations hinders monitoring schedule, cost, and other performance indicators, reducing the ability to reveal or manage the variability and uncertainty inherent in project activities (Alaloul et al. 2022). Because of these systemic bottlenecks, the industry persistently struggles to overcome the issues associated with safety hazards and low productivity (Jiang and Messner 2023). Consequently, there is a growing interest in the construction industry in adopting AI technologies for various management and construction tasks (Rabbi and Jeelani 2024), with computer vision emerging as a prominent solution to capture and analyze data from these complex sites (Paneru and Jeelani 2021).

To overcome the data deficiencies of traditional manual methods, research initially shifted toward automated object-level data collection (Tuomisto et al. 2026). Rather than broadly applying automation to all construction tasks, this specific technological domain focused on developing automated, real-time performance monitoring systems consisting of location tracking, activity recognition, and performance monitoring (Sherafat et al. 2020). To achieve this, the first of these methods involved attaching tracking devices like Radio Frequency Identification (RFID), Ultra-Wideband (UWB), or Bluetooth Low Energy (BLE) to building components (Tuomisto et al. 2026). However, while these tag-based methods are effective for locating materials, they cannot verify the proper installation or condition of components, limiting their utility for detailed inspection tasks (Tuomisto et al. 2026). Furthermore, tracking eventually thousands of items do not offer an economical or practical approach in large-scale projects (Bohn and Teizer 2010; Chen et al. 2023). Eventually, the industry shifted toward computer vision techniques to address this problem, as they operate without requirements to attach any sensors and enable more comprehensive on-site information, such as the locations and activities of construction entities, bringing operational and technical advantages over various wearable sensors (Xiong et al. 2019).

Building upon the shift toward camera-based monitoring, traditional computer vision approaches utilizing bounding-box detectors have been widely adopted to locate construction entities. While these models are highly effective at spatial localization, identifying only the workers and PPE components, without recognizing the contextual relationship between them, does not effectively verify PPE compliance (Nath, Behzadan, and Paal 2020). Specifically, inferring that workers are prone to experience falling from height based on safety regulations is difficult for computer vision; this condition can be summarized as a semantic gap (Wu et al. 2021). Because computer vision cannot natively comprehend safety rulebooks, traditional pipelines attempted to bridge this gap by applying an algorithm with explicit rules that are determined through empirical observations (Nath et al. 2020). However, manually establishing these rules for multiple PPEs is a complicated task and results may not be structured or comprehensive enough to consider all possible scenarios (Nath et al. 2020).

To overcome the limitation of manual geometric mapping, research explored an alternative approach using Natural Language Processing (NLP) for automatically extracting regulatory requirements from textual regulatory documents and formalizing these requirements in a logic format (Zhang and El-Gohary 2015). Yet, while NLP excels at rule extraction, these text-based models are fundamentally blind to the physical site. Ultimately, they simply represent the extracted safety requirements in the form of query graphs for supporting subsequent knowledge graph-based field compliance checking (Wang and El-Gohary 2023). Consequently, this creates a profound technological disconnect that computer vision can see the site but cannot understand the rules, while NLP understands the rules but cannot see the site. This fundamental limitation necessitates the emergence of multimodal Vision-Language Models (VLMs) capable of natively fusing both data streams.

To address the fundamental limitations of traditional methods in processing semantic information, recent advancements have introduced multi-modal large language models (MLLMs) upon which vision-language models (VLMs) combine image processing and natural language processing (NLP) (Chan et al. 2025). By extracting both intricate visual and linguistic features, these models can perform multiple types of instruction-following tasks such as image captioning, visual grounding, and visual question answering (VQA) (Chan et al. 2025). Furthermore, these architectures acquire generalized visual representations guided by natural language supervision, enabling impressive zero-shot detection capabilities (Zhang et al. 2025). By leveraging natural language prompts and contextual reasoning, VLMs offer a flexible approach for spatiotemporal analysis interpreting both the spatial geometry of hazards and the temporal sequence of complex human behaviors on dynamic construction sites, potentially reducing the dependence on extensive labeled datasets (Bui, Chodosh, and Tavakoli 2026). Consequently, these multimodal systems can create human-readable captions for daily progress and work activity logs (Jung et al. 2024), while simultaneously supporting interactive question-and-answer sessions, fostering a more informed understanding of construction site safety (Hussain et al. 2026).

Despite the rapid emergence of VLMs as a paradigm-shifting technology, existing systematic reviews have remained siloed in their respective unimodal domains. The majority of existing reviews on AI applications in the construction industry tend to focus on specific domains within AI, such as computer vision or NLP, thereby addressing only fragments of the potential AI applications (Rabbi and Jeelani 2024). For instance, previous reviews have extensively mapped pure computer vision research in construction (Jiang and Messner 2023; Paneru and Jeelani 2021; Zhong et al. 2019) or provided overviews of sensor-based safety management (Asadzadeh et al. 2020). While these studies provide valuable baselines, existing reviews do not cover the latest technological advances and lack systematic elaboration for construction scene understanding (Zhang et al. 2025). Specifically, although Vision-Language Models (VLMs) excel in general multimodal tasks, their ability to understand human behavior in safety-critical construction environments remains largely unexplored (Bui et al. 2026). Due to the development of scene understanding and large language models in recent years, it is necessary to conduct a new review on this topic to capture these paradigm-shifting advancements (Zhang et al. 2025).

To address this critical gap, this paper presents a comprehensive scoping review of multimodal vision-language systems applied to construction safety and progress monitoring. Conducted in strict accordance with the PRISMA-ScR guidelines (Tricco et al. 2018), the specific objectives of this study are to: (1) map the current landscape of VLM applications in construction management; (2) develop a robust taxonomy bridging foundational VLM architectures with practical construction management applications; and (3) analyze existing algorithmic limitations to outline a roadmap for future research in construction. To achieve these objectives, this review seeks to explicitly answer the following four primary research questions (RQs):

\begin{itemize}
\item RQ1: What are the prevalent Vision-Language Model (VLM) architectures currently utilized for spatiotemporal analysis in construction?
\item RQ2: How is multimodal data (visual and textual) fused and processed to evaluate safety hazards on active sites?
\item RQ3: What are the primary methodological bottlenecks in deploying these models for real-time progress tracking and safety?
\item RQ4: How is the performance of Vision-Language Models currently evaluated and benchmarked for spatiotemporal tasks in construction?
\end{itemize}
\section*{Research Methodology}
\coachcomment{Word comment at this heading}{Drafted according to PRISMA-ScR.}
This review was conducted in accordance with the PRISMA-ScR (Preferred Reporting Items for Systematic Reviews and Meta-Analyses extension for Scoping Reviews) guidelines (Tricco et al. 2018). A scoping review design was deliberately selected over traditional quantitative meta-analysis for two primary reasons. First, Vision-Language Models (VLMs) have emerged as a promising tool for visual understanding tasks (Bui et al. 2026), and recent studies are widely investigating these multi-modal feature mining pipelines to shed light on novel opportunities in the construction industry (Chan et al. 2025; Chen et al. 2024); thus, the primary objective is to map the extent, range, and nature of this novel architecture rather than to compare narrow intervention effects. Second, the evaluation metrics across the included studies exhibit substantial methodological heterogeneity; ranging from geometric precision metrics, such as mean Average Precision (Nath et al. 2020; Wang and El-Gohary 2024), to natural language generation metrics, such as BLEU and CIDEr (Jung et al. 2024; Xiao, Wang, and Kang 2022), which renders a standardized statistical meta-analysis inappropriate.

\subsection*{2.1 Search Strategy and Information Sources}
To ensure a rigorous and reproducible retrieval process, the literature search was conducted on 4th March 2026 across two of the most comprehensive scientific databases for engineering and computer science research: the Web of Science (WoS) Core Collection and Scopus. These specific databases were selected because the WoS contains the most comprehensive and influential journals within the field of construction and engineering management (Zhong et al. 2019), and Scopus is actively used alongside WoS as a primary search engine for AI applications in the construction domain (Rabbi and Jeelani 2024).

The keywords were determined by analyzing frequent words found in papers where AI was applied to enhance construction management and safety (Rabbi and Jeelani 2024). The search query was constructed using a tripartite Boolean logic framework designed to intersect three distinct conceptual domains: (1) the specific AI architecture (Multimodal vision-language systems), (2) the analytical task (spatiotemporal and dynamic reasoning), and (3) the application domain (construction management). The exact search strings applied across both databases are as follows:

String 1 - AI Architecture: ("Vision-Language Model*" OR "VLM" OR "Vision-Text" OR "Image-Text" OR "Visual Question Answering" OR "VQA" OR "Image Captioning" OR "CLIP" OR "Large Multimodal Model*" OR "LMM" OR "Foundation Model*" OR "Generative AI" OR "Video-Language" OR "Video-Text") AND

String 2 - Analytical Task: ("Spatiotemporal" OR "Spatio-temporal" OR "Video" OR "Temporal" OR "Dynamic" OR "Time-series" OR "4D" OR "Motion" OR "Tracking" OR "Sequenc*" OR "Action recognition" OR "Activity recognition") AND

String 3 - Application Domain: ("Construction Management" OR "Civil Engineering" OR "AEC" OR "Infrastructure" OR "Built Environment" OR "Construction Site" OR "Safety Monitoring" OR "Structural Health Monitoring" OR "Digital Twin")

Given the rapidly evolving nomenclature of VLM research, relying purely on an algorithmic, keyword-based search poses a risk of omitting foundational or transitional works that may not employ contemporary terminology. To address this limitation, a secondary backward snowballing (ancestry/reference chain) strategy was employed. The reference lists of all articles retained after the full-text review were systematically mined to identify foundational or transitional technologies, ensuring that some articles that were missed during the keyword search process were identified from using reference chains (Rabbi and Jeelani 2024).

\subsection*{2.2 Eligibility Criteria}
To maintain a disciplined and reproducible scope in accordance with the PRISMA-ScR guidelines (Tricco et al. 2018), explicit inclusion and exclusion criteria were established a priori, prior to the commencement of the screening process. These criteria were strategically designed to isolate high-quality literature that directly addresses the intersection of multimodal AI and active construction management, as narrowing the selection to reputable peer-reviewed journals enhances the reliability and relevance of the chosen studies (Rabbi and Jeelani 2024). Furthermore, the criteria deliberately retain a select number of traditional unimodal Computer Vision (CV) and Natural Language Processing (NLP) papers to establish an evolutionary comparative baseline. This baseline is essential to empirically demonstrate the inherent limitations of prior technologies and the "semantic gap," which indicates the limitation of computer vision systems in extracting the visual data of what the cameras "see" (Wu et al. 2021). The exact criteria applied during the screening process are summarized in Table 1.

Table 1. Summary of inclusion and exclusion criteria.

\tableblock{Dimension | Inclusion Criteria | Exclusion Criteria}{\textbf{Dimension:} Publication Date\\\\[0.25em]
\textbf{Inclusion Criteria:} 2019-2026 (for core VLM literature); 2010-2026 (for foundational CV baselines).\\\\[0.25em]
\textbf{Exclusion Criteria:} Pre-2010; post-March 2026.\par\medskip \textbf{Dimension:} Source Type\\\\[0.25em]
\textbf{Inclusion Criteria:} High-impact Peer-reviewed journal (Rabbi \& Jeelani, 2024).\\\\[0.25em]
\textbf{Exclusion Criteria:} Editorials, book chapters, theses, documents not written in English (Li, Deng, and Deng 2025), and conference proceedings and industry white papers (Rabbi and Jeelani 2024).\par\medskip \textbf{Dimension:} Technology\\\\[0.25em]
\textbf{Inclusion Criteria:} Multimodal architecture integrating visual/geometric data with textual/semantic data (e.g., VLMs, VQA, Image Captioning, CLIP-based models, MLLMs).\\\\[0.25em]
\textbf{Exclusion Criteria:} Purely unimodal CV (e.g., standalone YOLO, CNN bounding-box detectors) or purely text-based NLP, unless serving as an explicit comparative baseline within a hybrid framework.\par\medskip \textbf{Dimension:} Application Domain\\\\[0.25em]
\textbf{Inclusion Criteria:} Active construction phase: Safety Management (hazard identification, PPE compliance) and Progress Tracking (productivity analysis, schedule monitoring, reporting).\\\\[0.25em]
\textbf{Exclusion Criteria:} Post-construction, facility management, structural health monitoring, or non-engineering domains.}
\subsection*{2.3 Selection of Sources of Evidence}
The selection process for sources of evidence was conducted in two sequential phases: primary database retrieval and backward snowballing, and is visually summarized in the PRISMA-ScR flow diagram (Figure 1) (Tricco et al. 2018).

During Phase 1, the initial query across the Web of Science and Scopus databases yielded 1,443 raw records (769 from WoS and 674 from Scopus). Automated database filters restricting the results to relevant subject areas, the English language, and peer-reviewed document types excluded 1,320 records, leaving a targeted pool of 123 articles. Applying these filters is a standard practice to refine the scope of the results (Li et al. 2025) and to narrow the selection to articles published within reputable peer-reviewed journals, enhancing the reliability and relevance of the chosen studies (Rabbi and Jeelani 2024). After removing 33 duplicate records identified across the two databases, 90 unique articles remained for manual screening.

A rigorous two-stage manual screening protocol was subsequently employed to prevent selection bias, aligning with established scoping review methodologies (Tricco et al. 2018). In the first stage, title and abstract screening was conducted by a primary reviewer and validated by an expert academic reviewer, resulting in the exclusion of 32 articles that fell outside the scope of the active construction context. The remaining 58 articles were advanced to full-text review. During this second stage, 33 articles were excluded due to an exclusive focus on domains outside the active construction phase, insufficient methodological detail regarding their multimodal architectures, or reliance on unimodal architectures that did not belong to the baseline-to-VLM progression. This rigorous filtration process yielded a Phase 1 dataset of 25 core VLM and transitional articles.

Concurrently, Phase 2 was executed via a backward snowballing (reference list mining) search of the retained articles. This strategy was utilized because more relevant articles can be identified from these papers by using reference chains, which helps identify important articles that were missed during the keyword search process (Paneru and Jeelani 2021; Rabbi and Jeelani 2024). This ancestry search successfully identified an additional 31 highly relevant baseline and transition articles.

Integrating the Phase 1 database retrieval and Phase 2 snowball search culminated in a final, robust corpus of 56 articles included in this review. As represented in the final corpus synthesis, this dataset is methodologically balanced, comprising 31 baseline/background papers, 5 transition/bridging papers, and 20 core construction VLM papers. Curating the literature across this specific timeframe provides the necessary temporal data to present how these methods applied in the construction phase have evolved over the years (Jiang and Messner 2023).

\subsection*{2.4 Data Charting and Extraction Process}
Following the selection of the 56 core articles, a standardized data extraction framework was developed to systematically chart the literature. As recommended by scoping review methodologies (Tricco et al. 2018), the data charting form was iteratively refined by the research team to ensure all relevant variables addressing the research questions were captured consistently.

To provide a structured and comprehensive evaluation of the literature, the extracted data were systematically categorized into three distinct analytical dimensions: Content Analysis, Technical Analysis, and Critical Analysis. The complete data charting framework and its corresponding extraction variables are detailed in Table 2.

Table 2. Data Charting and Extraction Framework.

\tableblock{Analytical Dimension | Extraction Variables | Description}{\textbf{Analytical Dimension:} Dimension 1: Content Analysis\\\\[0.25em]
\textbf{Extraction Variables:} Cluster/Application; Input Data; Spatiotemporal Task; Automation Level.\\\\[0.25em]
\textbf{Description:} Identifies what the AI is doing by mapping the specific domain (e.g., safety vs. progress tracking), the sensory inputs utilized (e.g., video, image, point cloud), the specific spatiotemporal task (e.g., action recognition, tracking), and the system's level of autonomy (assistive vs. automated).\par\medskip \textbf{Analytical Dimension:} Dimension 2: Technical Analysis\\\\[0.25em]
\textbf{Extraction Variables:} Algorithm/Model; Dataset; Validation Metric; Compute/Hardware.\\\\[0.25em]
\textbf{Description:} Analyzes how the system operates by documenting the specific architectures (e.g., GPT-4o, CLIP), the nature of the training data (public vs. private/custom), the diverse evaluation metrics employed (e.g., mAP, F1-score), and the deployment environment (edge vs. cloud).\par\medskip \textbf{Analytical Dimension:} Dimension 3: Critical Analysis\\\\[0.25em]
\textbf{Extraction Variables:} Key Limitation; Research Gap; Future Direction.\\\\[0.25em]
\textbf{Description:} Evaluates the significance of the study by synthesizing the inherent technological or methodological bottlenecks, identifying unmet industry needs, and charting actionable future research trajectories.}
\subsection*{2.5 Data Synthesis}
Given the heterogeneous nature of the evaluation metrics across the included studies-which range from geometric intersection-over-union (IoU) and mean average precision (mAP) metrics for measuring bounding box predictions (Nath et al. 2020; Wang and El-Gohary 2024) to semantic fluency metrics for evaluating image captioning models, such as BLEU, METEOR, and CIDEr (Jung et al. 2024; Xiao et al. 2022); a standardized quantitative meta-analysis was deemed inappropriate. Instead, a thematic and narrative synthesis approach was adopted to map the literature (Tricco et al. 2018). To directly address the research questions regarding architectural evolution and multimodal data fusion, the extracted data were synthesized to present how these methods applied in the construction phase have evolved over the years (Jiang and Messner 2023). The final corpus of 56 papers was thematically classified into three distinct developmental stages: (1) the traditional baseline establishing the geometric-semantic disconnect, (2) the transitional bridge leveraging advanced but unimodal algorithms, and (3) the paradigm shift toward end-to-end Vision-Language Models (VLMs). This thematic categorization forms the structural foundation for the subsequent scientometric and thematic analyses presented in Chapters 3, 4, and 5.

\section*{Chapter 3 Planning Notes}
\coachcomment{Advisor methodology guidance}{Use descriptive mapping, not meta-analysis. The analysis should map publication trends, geography, study characteristics, topics/content, methods, and theoretical grounding.}
\begin{itemize}
\item Journal list, no of papers, citescore, SJR impact factor - table 2
\item Previous Review studies - table 1.
\item documents cocitation cluster table
\item highly cited articles
\item id-idf keyword for each year
\item top-10 cited authors,
\item Graphs - most frequently used keywords, frequency of technology used (heatmap), Network of frequently used keywords, document co citation networks, word cloud, co-occurences mapping of keywords, journal cocitation network, author cocitation map, collaboration network by institute, collaboration network of authors, Articles reviewed- categorized by application area, mind map
\end{itemize}
\coachcomment{Descriptive analysis note}{Frequently - not meta-analysis. Map: publication trends, geographic mapping, study characteristics (designs, methodologies, participants); research frequent methodologies and interventions.}
Number of studies, publication years, countries/regions, journals, types of publication

Study characteristics - research design, data sources, sample population, sample sizes, setting

Topic/Content characteristics: what outcomes examined, what was considered,

Methodological - modeling approaches, statistical techniques, ML methods, survey methods, logistic vs cross-sectional designs

Theories - Theories used, conceptual frameworks (or if there is not any theoretical grounding)


\coachcomment{Comment on removed Chapter 4 heading}{Jacqueline Borchert's Word comment at the original Chapter 4 heading was: ``Two major sections.'' This supports combining the findings/discussion structure rather than carrying the findings across separate Chapter 4 and Chapter 5 sections. The original Chapter 4, Chapter 5, and bibliography are omitted from this PDF.}
\end{document}

